% !TEX root = ../main.tex

\begin{acknowledgements}
时光荏苒，本科四年的学习生涯即将画上句号。回顾毕业设计的研究历程，我深感收获良多，在此向所有给予我帮助和支持的人表达最诚挚的感谢。

首先，我要衷心感谢我的指导老师李成林老师。从选题方向的确定到研究方案的论证，从实验设计的推敲到论文写作的修改，李老师始终给予我耐心细致的指导与严格的学术要求。李老师严谨的治学态度和开阔的学术视野，不仅帮助我顺利完成了本课题的研究，更让我在科研方法论上受益匪浅。

其次，我要特别感谢博士生学长阚诺文。在研究过程中，阚学长在强化学习算法选型、实验方案设计以及工程实现等多个环节给予了我宝贵的指导意见和技术建议。每当我在研究中遇到困难和瓶颈时，阚学长总是不厌其烦地与我讨论分析，帮助我理清思路、找到突破口。他丰富的科研经验和乐于助人的品格令我深受感动。

我还要感谢我实习所在的公司字节跳动。公司允许我在工作期间抽出时间进行毕业设计的研究工作，为我提供了灵活的时间安排，使我能够兼顾实习与学业。同时，公司内部公开的大量实时通信、网络传输与机器学习相关的技术资料和学习课程，极大地拓宽了我的技术视野，为本课题的研究提供了重要的知识储备与启发。

最后，我要感谢我的父母。感谢他们多年来对我学业的无条件支持与默默付出，无论是在学习上遇到挫折还是面临人生选择时，他们始终是我最坚实的后盾。正是他们的理解、鼓励与关爱，让我能够心无旁骛地投入学业和研究之中。

谨以此文献给所有关心和帮助过我的人。
\end{acknowledgements}
